\documentclass[a4paper]{article}
\usepackage{bookmark}
\usepackage[italian]{babel}
\usepackage{geometry}
\usepackage{hyperref}
\usepackage{sectsty}
\usepackage{parskip}

\hypersetup{hidelinks}
\allsectionsfont{\centering}

\begin{document}
    
    \title{Analisi dell'amplificatore Miller}
    \author{Luca Brescia}
    \date{\today}
    \maketitle

    \section*{Obiettivo}
        Creazione di uno standard per le simulazioni dei circuiti della tesi. \\
        Utilizzare prima il circuito di un amplificatore Miller, in configurazione di inseguitore (buffer), a cui vanno effettuate le seguenti simulazioni:

        \begin{itemize}
            \item Analis DC: in cui bisogna osservare l'assestamento dell'uscita al valore prefissato di ingresso.
            \item Analisi AC: fare sweep in frequenza da 1 MHz a 1-2-4 GHz
        \end{itemize}

        L'ingresso dovrà essere impostato su 0.5Vpp e ampiezza 1Vpp basato su una tensione di alimentazione di Vdd=1V.

        Inoltre si dovranno effettuare le simulazioni con le seguenti configurazioni:

        \begin{itemize}
            \item 1 stadio;
            \item 2 stadi;
            \item Eff. corner;
            \item PVT
            \item Montecarlo;
        \end{itemize}

        Per gli ultimi tre punti riferirsi alla prof.ssa Richelli.
    \section*{Svolgimento}

\end{document}